% Teil II: der durchgerechnete Weg. Reihenfolge beibehalten.

\subsection*{Gegeben und gesucht}
Einige Funktionen zum Üben der Ableitungsregeln (Teil a) sowie Kurvendiskussion
von $f(x)=\ln(x^2+1)$ (Teil b). Einzelne konkrete Funktionen in Teil a sind im
eingescannten Blatt unlesbar; die Methode zeigt der lesbare Teil.

\subsection*{a) Ableiten — lesbarer Teil}
Mit der Kettenregel, z. B. für $g(x)=\ln(\cos x)$ mit $x\in(-\tfrac\pi2,\tfrac\pi2)$:
\begin{equation}
  g'(x)=\frac{1}{\cos x}\cdot(-\sin x)=-\tan x.
\end{equation}
Für einen Term wie $x^2\ln x$ von Typ Produkt hilft die Produktregel, für einen
Quotienten wie $\frac{1}{1+x^2}$ die Quotienten- bzw.\ Potenz-Kettenregel:
\begin{equation}
  \dv{}{x}\frac{1}{1+x^2}=-\frac{2x}{(1+x^2)^2}.
\end{equation}
Die restlichen Funktionen (deren genaue Gestalt der Scan nicht sicher zeigt)
ebenso mit den passenden Regeln ableiten — bitte die ursprünglichen Ausdrücke
vom Blatt übernehmen.

\subsection*{b) Kurvendiskussion von $f(x)=\ln(x^2+1)$}
Definitionsbereich: $x^2+1>0$ für alle $x$, also $\mathbb D_f=\mathbb R$.
Nullstellen:
\begin{equation}
  \ln(x^2+1)=0 \;\Longleftrightarrow\; x^2+1=1 \;\Longleftrightarrow\; x=0.
\end{equation}
Ableitungen:
\begin{equation}
  f'(x)=\frac{2x}{x^2+1},\qquad
  f''(x)=\frac{2(x^2+1)-2x\cdot2x}{(x^2+1)^2}
  =\frac{2-2x^2}{(x^2+1)^2}.
\end{equation}
Extrema: $f'=0$ bei $x=0$; dort $f''(0)=2>0$, also lokales Minimum $f(0)=0$,
zugleich globales Minimum. Wendepunkte: $f''=0$ bei $2-2x^2=0$, also
$x=\pm1$; dort wechselt $f''$ das Vorzeichen, daher sind $\pm1$ Wendestellen.

\begin{ergebnis}
\begin{equation*}
  f\ \text{überall definiert},\quad
  \text{Min}\ (0,0),\quad
  \text{Wende st.}\ x=\pm1.
\end{equation*}
\end{ergebnis}

\begin{kontrolle}
\textbf{Symmetrie:} $f(x)=f(-x)$ (gerade), also Min bei $0$ und symmetrische
Wendestellen bei $\pm1$ — plausibel.
\textbf{Grenzfall:} $f''$ ist für $|x|<1$ positiv, für $|x|>1$ negativ;
Vorzeichenwechsel bei $\pm1$ bestätigt die Wendestellen.
\end{kontrolle}

\begin{merksatz}
Kettenregel zuerst; bei $\ln(x^2+1)$ ist $x^2+1$ stets positiv, daher keine
Definitionslücken. Extrema über $f'=0$ mit Vorzeichen, Wendepunkte über $f''=0$.
\end{merksatz}
